% Source PDF: source/普通高考/2016/2016全国3理(云南,广西,贵州).pdf, page 1.
% PDF embedded-bitmap attempt: pdfimages -list found no embedded image; the flowchart is vector line art.
% Grid evidence: tmp/review_flow_2016_np3_science/page-1_grid100.jpg and q7_block_grid50.jpg;
% comparison source is cropped from the ungridded page render (q7_original.png).
% Elements: rounded start/end terminals, input/output parallelograms, process boxes,
% decision diamond, directed arrows, and 是/否 branch labels.
% Node -> directed edge list: 开始→输入 a,b→(n=0,s=0)→(a=b−a)→(b=b−a)→(a=b+a)
% →(s=s+a,n=n+1)→(s>16); 是→输出 n→结束; 否→(left loop)→(a=b−a).
% solid segments: all box outlines and connector lines/arrows; dashed segments: none.
% labels: all node text is locally zero-indented and centered; 否 sits by the left loop,
% 是 sits beside the downward true branch.  Loop and output paths are disjoint.

\begin{tikzpicture}[x=1cm,y=0.90cm,>=Stealth,
  line width=0.45pt,line cap=round,line join=round,
  every node/.style={font=\normalsize,anchor=center,align=center,
    execute at begin node={\setlength{\parindent}{0pt}}},
  flowbox/.style={draw,minimum height=0.62cm,inner sep=0pt,align=center},
  io/.style={draw,shape=trapezium,trapezium left angle=78,trapezium right angle=102,
    trapezium stretches=false,
    minimum height=0.72cm,inner sep=0pt,align=center,xscale=0.70},
  flowarrow/.style={->,line width=0.45pt},
]
  % Main vertical sequence.
  \node[draw,rounded corners=0.16cm,minimum width=1.10cm,minimum height=0.62cm] (start) at (0,10.80) {开始};
  \node[io,minimum width=0pt] (input) at (0,9.58) {输入 $a,b$};
  \node[flowbox,minimum width=2.30cm] (init) at (0,8.34) {$n=0,\ s=0$};
  \node[flowbox,minimum width=1.75cm] (aone) at (0,7.10) {$a=b-a$};
  \node[flowbox,minimum width=1.75cm] (bone) at (0,5.86) {$b=b-a$};
  \node[flowbox,minimum width=1.75cm] (atwo) at (0,4.62) {$a=b+a$};
  \node[flowbox,minimum width=3.60cm] (update) at (0,3.38) {$s=s+a,\ n=n+1$};
  \node[draw,diamond,aspect=2.75,minimum width=2.72cm,minimum height=0.84cm,inner sep=0pt,align=center] (test) at (0,2.05) {$s>16$};
  \node[io,minimum width=0pt] (output) at (0,0.78) {输出 $n$};
  \node[draw,rounded corners=0.16cm,minimum width=1.10cm,minimum height=0.62cm] (finish) at (0,-0.48) {结束};

  \draw[flowarrow] (start.south) -- (input.north);
  \draw[flowarrow] (input.south) -- (init.north);
  % The loop joins the trunk above aone; only this single downward arrow enters it.
  \coordinate (loop-top) at (0,7.47);
  \draw (init.south) -- (loop-top);
  \draw[flowarrow] (loop-top) -- (aone.north);
  \draw[flowarrow] (aone.south) -- (bone.north);
  \draw[flowarrow] (bone.south) -- (atwo.north);
  \draw[flowarrow] (atwo.south) -- (update.north);
  \draw[flowarrow] (update.south) -- (test.north);
  \draw[flowarrow] (test.south) -- (output.north);
  \draw[flowarrow] (output.south) -- (finish.north);

  % 否 loop: route left, then rejoin the main trunk before its single downward arrow into a=b-a.
  \coordinate (left) at (-2.28,2.05);
  \coordinate (entry) at (-2.28,7.47);
  \draw (test.west) -- (left) -- (entry);
  \draw (entry) -- (loop-top);

  \node[anchor=east,inner sep=0pt] at (-1.58,2.52) {否};
  \node[anchor=west,inner sep=0pt] at (0.20,1.30) {是};

  % Keep a restrained crop margin like the source page crop; no text is offset by hand.
  \path[use as bounding box] (-3.05,11.15) rectangle (2.65,-0.82);
\end{tikzpicture}
